\begin{lemma}\label{push:src:A.3.1}\label{push:A31}
Let $Z$ be a compact reduced local-complete-intersection complex
space of pure dimension $n$. In any proper flat holomorphic
family over a disc with central fibre $Z$ and smooth nearby
fibres, their Euler number is
\[
 e_{\mathrm{vir}}(Z)=\langle c_n(T_Z^{\mathrm{vir}}),[Z]\rangle.
\]
This statement needs no projectivity or K\"ahler hypothesis.
\end{lemma}
\Uses{the hypotheses, notation and definitions stated here; \cite[Lemmas 3.1--3.2 and Section 3.1]{Levy08}; \cite[Section~3, Theorems~1--2, pp.~463--464]{Loj64}; \cite[Theorem~1.6, p.~5; Remark~1.16, p.~11]{BM95AB}; \cite{FG10}}

\begin{proof}
Flat deformation preserves the local-complete-intersection
property relative to the disc near $Z$. Locally this follows by
lifting an embedding and its regular sequence of equations;
Nakayama's lemma and flatness show that their lifts generate
the ideal and remain a regular sequence. Properness allows one
to shrink the disc so that these neighbourhoods cover the total
space. The relative cotangent complex is consequently perfect
of amplitude $[-1,0]$. Its dual defines a topological complex
$K$-class $V$ on that total space; local finite vector-bundle
complexes define the class by descent, or by stabilization and
partitions of unity; the analytic-to-topological comparison is
given in \cite[Lemmas 3.1--3.2 and Section 3.1]{Levy08}.
Its restriction to a smooth fibre is its
tangent bundle, and its restriction to $Z$ is its intrinsic
virtual tangent class.

We compare the fundamental cycles on a smooth proper modification,
so that the ambient-manifold hypothesis of the triangulation theorem
is satisfied. Analytic resolution, in the complex-analytic scope of
\cite[Section~0.2(2), p.~3; Theorem~1.6 and its noncompact extension,
p.~5; Remark~1.16, p.~11]{BM95AB}, gives
\[
 \rho:\widetilde{\mathcal X}\longrightarrow\mathcal X
\]
with smooth total space, isomorphic over the smooth locus. Only a
neighbourhood of the compact central fibre is needed; the locally
finite construction is finite there after shrinking the base.
The total space is reduced: a nilpotent local section lies in every
power of the base parameter by reducedness of the central fibre and
flatness, and hence vanishes by Krull intersection. Thus the stated
resolution applies to the actual dominating space.

Write the divisor of the pulled-back parameter as
$\sum_j m_jD_j$. The inverse image of a short radial segment
$[0,\epsilon]$ is a compact semianalytic subset of the smooth real
analytic manifold $\widetilde{\mathcal X}$. We triangulate it,
compatibly with its two endpoint fibres and the singular loci of the
$D_j$, by \cite[Section~3, Theorems~1--2, pp.~463--464]{Loj64}.
Over the interior of the segment its top-dimensional part is an
oriented real $(2n+1)$-manifold. Near a general point of $D_j$ the
parameter is $u z^{m_j}$, with $u$ a holomorphic unit; choosing a
local root of $u$ exhibits $m_j$ radial half-sheets, each with the
complex boundary orientation. The sum of their oriented top
simplices therefore has boundary
\[
 [\mathcal X_\epsilon]-\sum_jm_j[D_j].
\]
The singular and crossing sets have real dimension at most $2n-2$
and contribute no boundary simplex of dimension $2n$.

At every smooth point of the reduced central fibre the original
map is a submersion, by flatness and smoothness of that fibre germ.
The modification is consequently an isomorphism there and the
strict-transform multiplicity is one. Every other $D_j$ maps into
$\operatorname{Sing}Z$, of complex dimension at most $n-1$.
Its fundamental class has zero pushforward in dimension $2n$:
its map factors through that lower-dimensional analytic set.
Thus pushing the displayed chain forward gives exactly
$[\mathcal X_\epsilon]-[Z]$. This proves the required homology
identity in the original total space. Pairing the globally defined
$c_n(V)$ with these homologous cycles gives the equality of
Chern numbers. Gauss--Bonnet on the smooth compact complex
fibre gives its Euler number. This proves the analytic assertion
directly. The algebraic counterpart is Theorem 5.8 and
Remark 5.9 of \cite{FG10}; those scheme statements alone would
not justify omitting the analytic comparison above.
\end{proof}